\section{การออกแบบและพัฒนาระบบ}
\subsection{สถาปัตยกรรมของระบบ}
\begin{figure}[H]
\centering
\begin{tikzpicture}[
  x=1mm,y=1mm,
  box/.style={draw,rectangle,align=center,text width=32mm,minimum height=10mm,inner sep=1.5mm,font=\fontsize{12bp}{16bp}\selectfont},
  flow/.style={-{Stealth[length=1.7mm]},line width=0.5pt},
  exchange/.style={{Stealth[length=1.7mm]}-{Stealth[length=1.7mm]},line width=0.5pt}
]
\node[box] (web) at (0,0) {ผู้ใช้ผ่านเว็บระบบใบลา};
\node[box] (api) at (0,-18) {บริการระบบใบลา};
\node[box] (support) at (43,-18) {ระบบข้อมูลสนับสนุน\\วิชา ผู้สอน ที่ปรึกษา};
\node[box] (db) at (0,-39) {ฐานข้อมูลใบลา\\และ snapshot};
\node[box] (etl) at (43,-39) {ETL\\ตรวจและแปลงข้อมูล};
\node[box] (report) at (0,-60) {บริการรายงาน\\ตรวจสิทธิ์ก่อนส่งผล};
\node[box] (dw) at (43,-60) {คลังข้อมูลรายงาน};
\draw[exchange] (web) -- (api);
\draw[exchange] (api) -- (support);
\draw[exchange] (api) -- (db);
\draw[flow] (db) -- (etl);
\draw[flow] (support) -- (etl);
\draw[flow] (etl) -- (dw);
\draw[exchange] (report) -- (dw);
\draw[exchange] (report.west) -- (-20,-60) -- (-20,0) -- (web.west);
\end{tikzpicture}
\caption{องค์ประกอบและการเชื่อมต่อข้อมูลของระบบ}\label{fig:architecture}
\end{figure}
ส่วนติดต่อผู้ใช้พัฒนาด้วย Next.js และ React ส่วนบริการประมวลผลใช้ NestJS และเข้าถึงฐานข้อมูลผ่าน Prisma ตามแฟ้มกำหนดแพ็กเกจและสคีมาในโครงการ ระบบเชื่อมบริการยืนยันตัวตน บุคลากร และข้อมูลประกอบใบลาจากภายนอก โดยการตั้งค่าปลายทางและข้อมูลลับอยู่ฝั่งเซิร์ฟเวอร์ แผนภาพสรุปองค์ประกอบหลักอยู่ในภาพที่~\ref{fig:architecture} ลูกศรสองทางหมายถึงการเรียกบริการและรับผล ส่วนลูกศรทางเดียวหมายถึงการส่งข้อมูลเข้าสู่กระบวนการรายงาน โดยบริการรายงานตรวจขอบเขตผู้ใช้กับต้นทางก่อนส่งผลให้หน้าเว็บ

\subsubsection{เทคโนโลยีส่วนติดต่อผู้ใช้}
เว็บระบบใบลาและระบบข้อมูลสนับสนุนใช้ Next.js และ React สำหรับสร้างหน้าเว็บร่วมกับ TypeScript เพื่อจัดการโครงสร้างข้อมูล ส่วนติดต่อแบ่งเป็นหน้าสำหรับนิสิต อาจารย์ และผู้ดูแลระบบตามหน้าที่ของแต่ละกลุ่ม \cite{repo}
ส่วนติดต่อใช้ Tailwind CSS 4 ร่วมกับคอมโพเนนต์แนวทาง shadcn/ui และ Radix UI เพื่อประกอบแบบฟอร์ม ตาราง และกล่องโต้ตอบ การแยกคอมโพเนนต์ช่วยให้หน้าสำหรับนิสิต อาจารย์ และผู้ดูแลระบบใช้รูปแบบร่วมกันได้ ขณะที่ตรรกะเฉพาะบทบาทยังแยกในหน้าหรือบริการที่เกี่ยวข้อง
แบบฟอร์มยื่นใบลาใช้ React Hook Form จัดการค่าช่องกรอกและสถานะฟอร์ม ร่วมกับ Zod และ resolver สำหรับตรวจรูปแบบข้อมูลก่อนส่ง การตรวจฝั่งหน้าเว็บช่วยให้แจ้งข้อผิดพลาดใกล้ช่องกรอก แต่ไม่แทนการตรวจเงื่อนไขและสิทธิ์ฝั่งเซิร์ฟเวอร์ Axios ใช้รับส่งข้อมูลกับบริการ ส่วน Recharts ใช้สร้างกราฟรายงาน และส่วนเอกสารใช้ html-to-image ร่วมกับ jsPDF ในขั้นตอนส่งออกเอกสาร
\subsubsection{เทคโนโลยีบริการและการจัดเก็บ}
บริการทั้งสองโครงการใช้ NestJS บน Node.js โดยแยกส่วนรับคำขอ กฎการทำงาน และความสามารถแต่ละด้าน เช่น ใบลา ข้อมูลสนับสนุน และรายงาน โครงสร้างนี้ทำให้ตรรกะการเปลี่ยนสถานะไม่จำเป็นต้องอยู่ในหน้าเว็บ และใช้ซ้ำได้จากผู้เรียกที่ได้รับสิทธิ์
Prisma เชื่อมต่อฐานข้อมูล MySQL เพื่อจัดเก็บและค้นหาข้อมูลที่สัมพันธ์กัน ส่วนการเขียนคำขอกับขั้นอนุมัติใช้ธุรกรรมเพื่อรักษาความสอดคล้อง Multer รองรับการรับแฟ้มผ่านบริการฝั่งเซิร์ฟเวอร์ โดยการตรวจเจ้าของและสถานะคำขอเป็นตรรกะอีกส่วนหนึ่งของระบบใบลา
ตารางที่~\ref{tab:tech-stack} สรุปเทคโนโลยีตามหน้าที่ เพื่อแยกเฟรมเวิร์กส่วนติดต่อ เครื่องมือจัดเก็บ และเครื่องมือวิเคราะห์ออกจากกัน รายการแพ็กเกจที่มีอยู่เพียงอย่างเดียวไม่ใช่หลักฐานว่าทุกความสามารถถูกเปิดใช้ในการติดตั้งจริง
\begin{table}[H]
\caption{เทคโนโลยีหลักและหน้าที่ในระบบ}\label{tab:tech-stack}
\centering
\begin{tabularx}{\columnwidth}{|p{30mm}|X|}
\hline
เทคโนโลยี & หน้าที่ \\ \hline
Next.js / React & หน้าเว็บและคอมโพเนนต์ผู้ใช้ \\ \hline
Tailwind / shadcn/ui & รูปแบบและองค์ประกอบหน้าจอ \\ \hline
React Hook Form / Zod & จัดการและตรวจข้อมูลแบบฟอร์ม \\ \hline
NestJS / TypeScript & บริการและกฎการทำงาน \\ \hline
Prisma / MySQL & เชื่อมข้อมูลและจัดเก็บธุรกรรม \\ \hline
Recharts & กราฟสถิติการลา \\ \hline
xlsx & อ่านข้อมูล Excel ในระบบสนับสนุน \\ \hline
Node.js / SQL & ประมวลผล ETL และคลังข้อมูล \\ \hline
\end{tabularx}
\end{table}
\subsubsection{เทคโนโลยีข้อมูลสนับสนุนและคลังข้อมูล}
บริการนำเข้าของระบบข้อมูลสนับสนุนใช้ไลบรารี xlsx อ่าน workbook และแปลงแถวเป็นข้อมูลวิชา ตอนเรียน ผู้สอน หรือที่ปรึกษา ก่อนเขียนด้วย Prisma การจัดการรูปแบบ Excel จึงเป็นหน้าที่ของบริการนำเข้า ส่วนกฎรายการซ้ำและความสัมพันธ์อาศัยทั้งตรรกะโปรแกรมและคีย์ของฐานข้อมูล
คลังข้อมูลใช้ Node.js จัดเตรียมข้อมูลสำหรับรายงาน โดยแยกการรวบรวม ตรวจสอบ และนำเข้าข้อมูลออกจากงานยื่นใบลา ช่วยให้กำหนดวิธีรวมยอดและปรับปรุงรายงานได้อย่างเป็นระบบ อีกทั้งลดการคำนวณรายงานซ้ำจากข้อมูลที่กำลังใช้พิจารณาคำขอ
\subsubsection{เครื่องมือพัฒนาและการติดตั้ง}
โครงการจัดเป็น pnpm monorepo แยก apps/web และ apps/api และมี lockfile ควบคุมการติดตั้ง dependency ระบบใบลาระบุ Node.js 22 และ pnpm 10.0.0 ในคู่มือและ Dockerfile โดยมี image สำหรับเว็บและบริการแยกกัน การแยกส่วนนี้ทำให้กำหนดตัวแปรแวดล้อมและพื้นที่เก็บแฟ้มของแต่ละบริการได้ตามหน้าที่
ESLint และ TypeScript ใช้ตรวจซอร์สโค้ด ส่วนชุดทดสอบสิทธิ์และ ETL ของระบบใบลาใช้ตัวทดสอบของ Node.js สำหรับระบบข้อมูลสนับสนุนมีการกำหนด Jest และ Supertest ในโครงการ การกล่าวถึงเครื่องมือเหล่านี้เป็นคำอธิบายสิ่งที่มีในโค้ด ไม่ใช่รายงานว่ารันผ่านทุกชุดในการจัดทำบทความครั้งนี้
\subsection{ระบบข้อมูลสนับสนุน}
ระบบข้อมูลสนับสนุนการลานิสิต คณะวิศวกรรมศาสตร์ กำแพงแสน (Student Leave Reference Data Management System) จัดการข้อมูลรายวิชา ตอนเรียน ผู้สอน และความสัมพันธ์ระหว่างที่ปรึกษากับนิสิตแยกจากข้อมูลคำขอลา ระบบใบลาเรียกใช้ข้อมูลเหล่านี้เพื่อประกอบแบบฟอร์ม เลือกวิชาและผู้สอน ตลอดจนตรวจขอบเขตนิสิตในความดูแล \cite{components}
ข้อมูลรายวิชารองรับหลายตอนเรียนและผู้สอนร่วม มีการเก็บตารางเรียน วันเรียน และช่วงเวลาเพื่อระบุว่าการลาส่งผลต่อการเรียนวิชาใดบ้าง ส่วนข้อมูลที่ปรึกษาช่วยกำหนดผู้พิจารณาขั้นแรกและขอบเขตรายงาน ความสัมพันธ์เหล่านี้จึงใช้ทั้งในการยื่นคำขอและการเข้าถึงข้อมูลตามหน้าที่
ผู้ดูแลข้อมูลสามารถนำเข้ารายวิชาและข้อมูลที่ปรึกษาจาก Excel โดยเลือกเพิ่มรายการใหม่หรือแทนที่ชุดเดิม ระบบตรวจรายการซ้ำและความสัมพันธ์ก่อนบันทึก การแยกงานนำเข้าออกจากการยื่นใบลาช่วยให้ปรับข้อมูลอ้างอิงได้เป็นรอบ และลดการกรอกข้อมูลเดียวกันซ้ำจากผู้ใช้หลายคน
ระบบใบลาเก็บตารางเรียนตามภาคการศึกษาและรายการวิชาที่นิสิตบันทึกไว้เพื่อเลือกใช้ซ้ำ ขณะที่ระบบข้อมูลสนับสนุนให้ข้อมูลอ้างอิงปัจจุบัน ประวัติใบลาจึงเก็บรายละเอียด ณ เวลายื่นแยกไว้ เพื่อให้การเปลี่ยนตารางเรียน ผู้สอน หรือที่ปรึกษาภายหลังไม่เปลี่ยนความหมายของใบลาเดิม
\subsection{การยื่นคำขอและการจัดเก็บข้อมูล}
นิสิตเลือกประเภทลา ระบุช่วงวันที่ เหตุผล รายวิชาที่ได้รับผลกระทบ และแนบหลักฐานตามเงื่อนไข ระบบเชื่อมคำขอกับข้อมูลนิสิตและภาคการศึกษา พร้อมเก็บสำเนารายละเอียด ณ เวลายื่น ได้แก่ ข้อมูลติดต่อ ที่ปรึกษา รายวิชา ตอนเรียน ผู้สอน และวันเรียนที่ลา เพื่อใช้แสดงเอกสารและตรวจสอบย้อนหลัง \cite{schema}
เส้นทางพิจารณาเริ่มจากอาจารย์ที่ปรึกษาแล้วส่งต่อหัวหน้าภาควิชา โดยเพิ่มรองคณบดีเมื่อจำนวนวันลาเกิน 15 วัน ดังภาพที่~\ref{fig:approval-route} ระบบนับวันตามปฏิทินทำการและวันเรียนพิเศษที่ได้รับผลกระทบ วันหยุดจึงไม่นับโดยอัตโนมัติ แต่สามารถนับได้เมื่อมีการเรียนในวันนั้น โดยวันเดียวกันไม่ถูกนับซ้ำแม้กระทบหลายวิชา
ระบบบันทึกคำขอกับลำดับผู้พิจารณาร่วมกัน เพื่อให้ติดตามความคืบหน้าได้ตลอดกระบวนการ ส่วนสำเนารายละเอียดช่วยรักษาบริบทเดิมแม้ข้อมูลอ้างอิงเปลี่ยน การสำรองข้อมูลจึงต้องครอบคลุมทั้งฐานข้อมูลและแฟ้มเอกสารที่คำขออ้างถึง
\subsection{การพิจารณาและการรับทราบ}
ผู้พิจารณาตัดสินได้เฉพาะคำขอที่ถึงขั้นของตน เมื่ออนุมัติระบบส่งต่อขั้นถัดไป และเมื่อครบทุกขั้นจึงเปลี่ยนใบลาเป็นอนุมัติ หากขั้นใดปฏิเสธ กระบวนการพิจารณาสิ้นสุดลง ระบบควบคุมการบันทึกคำตัดสินเพื่อป้องกันการส่งผลซ้ำหรือการตัดสินพร้อมกันเขียนทับผลเดิม
ผู้พิจารณาสามารถตั้งกฎอนุมัติอัตโนมัติสำหรับขั้นที่รับผิดชอบได้ โดยเริ่มนับเวลาทำการตั้งแต่คำขอเข้าสู่ขั้นนั้น และอนุมัติเมื่อครบระยะเวลาที่กำหนด ระบบแยกประวัติการตัดสินโดยบุคคลออกจากกฎอัตโนมัติเพื่อให้ตรวจสอบย้อนหลังได้ การใช้กฎดังกล่าวต้องสอดคล้องกับแนวปฏิบัติของหน่วยงาน
การพิจารณาแทนในระดับภาควิชาหรือคณะต้องตรวจสิทธิ์และขอบเขตหน่วยงาน พร้อมบันทึกทั้งผู้รับผิดชอบเดิมและผู้ตัดสินจริง ผู้ดูแลระบบจึงไม่ได้รับสิทธิ์อนุมัติแทนโดยอัตโนมัติ ส่วนอาจารย์ผู้สอนรับทราบเฉพาะใบลาที่อนุมัติแล้วและเกี่ยวข้องกับวิชาที่สอน การรับทราบไม่เปลี่ยนผลอนุมัติและไม่นับเป็นขั้นพิจารณาเพิ่มเติม
\subsection{แบบจำลองข้อมูลธุรกรรม}
ข้อมูลธุรกรรมแบ่งเป็นบัญชีผู้ใช้ นิสิต ใบลา ขั้นพิจารณา การรับทราบ และข้อมูลประกอบการศึกษา โดยใบลาหนึ่งใบมีผู้พิจารณาได้หลายขั้น แต่ละขั้นเก็บผู้รับผิดชอบ ลำดับ ผลตัดสิน เหตุผล และเวลาแยกกัน เพื่อให้ทราบทั้งสถานะรวมของคำขอและความคืบหน้ารายขั้น ตารางที่~\ref{tab:operational-model} สรุปหน้าที่ของข้อมูลแต่ละกลุ่ม
\begin{table}[H]
\caption{กลุ่มข้อมูลหลักของฐานธุรกรรม}\label{tab:operational-model}
\centering
\begin{tabularx}{\columnwidth}{|p{23mm}|X|}
\hline
กลุ่มข้อมูล & บทบาทในกระบวนการ \\ \hline
บัญชีและนิสิต & เชื่อมตัวตนผู้ยื่นกับคำขอและประวัติ \\ \hline
ใบลา & ประเภท ช่วงวันที่ จำนวนวัน สถานะ และตำแหน่งแฟ้ม \\ \hline
ขั้นพิจารณา & ผู้รับผิดชอบ บทบาท ลำดับ ผลตัดสิน และเวลา \\ \hline
การรับทราบ & เชื่อมใบลา วิชา ตอนเรียน และผู้สอน \\ \hline
ภาคการศึกษา & กำหนดบริบทปีและเทอมของคำขอ \\ \hline
\end{tabularx}
\end{table}
ระบบเก็บข้อมูลรายวิชา ตารางเรียน วันเรียน และผู้สอนเพื่อเชื่อมการลากับการเรียนที่ได้รับผลกระทบ การแยกข้อมูลอ้างอิงออกจากประวัติคำขอช่วยให้ปรับปรุงข้อมูลปัจจุบันได้โดยยังคงรายละเอียดเดิมสำหรับตรวจสอบย้อนหลัง แผนภาพความสัมพันธ์ของข้อมูลแสดงในภาคผนวก~\ref{app:db}
\subsection{สถานะคำขอและเงื่อนไขการเปลี่ยนขั้น}
สถานะคำขอแสดงว่ากระบวนการเสร็จสิ้นหรือยัง ส่วนสถานะขั้นพิจารณาแสดงว่าขณะนั้นอยู่ในความรับผิดชอบของใคร ใบลาจึงยังรอผลแม้ที่ปรึกษาอนุมัติแล้ว หากยังมีขั้นถัดไป ผู้ใช้งานเห็นความคืบหน้าตามลำดับดังตารางที่~\ref{tab:state-transition}
\begin{table}[H]
\caption{เหตุการณ์และการเปลี่ยนสถานะคำขอ}\label{tab:state-transition}
\centering
\begin{tabularx}{\columnwidth}{|p{25mm}|X|}
\hline
เหตุการณ์ & ผลต่อกระบวนการ \\ \hline
สร้างคำขอ & ใบลา รอพิจารณา; เปิดขั้นแรก \\ \hline
อนุมัติระหว่างทาง & บันทึกผลขั้นเดิมและเปิดขั้นถัดไป \\ \hline
อนุมัติขั้นสุดท้าย & ใบลาได้รับอนุมัติ \\ \hline
ปฏิเสธขั้นปัจจุบัน & ใบลาถูกปฏิเสธ \\ \hline
ส่งผลตัดสินซ้ำ & ต้องไม่เขียนทับขั้นที่ตัดสินเสร็จแล้ว \\ \hline
\end{tabularx}
\end{table}
ก่อนเปลี่ยนสถานะ ระบบตรวจตัวตน บทบาท ความสัมพันธ์กับคำขอ และขั้นที่เปิดให้พิจารณาร่วมกัน การตรวจเฉพาะชื่อบทบาทจึงไม่เพียงพอ เพราะผู้มีบทบาทเดียวกันอาจรับผิดชอบนิสิตหรือหน่วยงานต่างกัน
\subsection{การควบคุมหลักฐานและเอกสาร}
เจ้าของคำขออัปโหลดหรือแก้ไขหลักฐานได้ตามสถานะที่ระบบอนุญาต ผู้เข้าถึงเอกสารแนบต้องผ่านการตรวจสิทธิ์ของใบลานั้น ส่วนเอกสารใบลามี QR สำหรับเปิดหน้าตรวจข้อมูลและลำดับการพิจารณา ลิงก์ดังกล่าวแยกจากการเข้าถึงหลักฐานส่วนตัว เพื่อให้ตรวจเอกสารได้โดยไม่เผยแพร่สิ่งแนบทั้งหมด
\subsection{รายงานตามบทบาท}
นิสิตดูประวัติของตน ที่ปรึกษาดูนิสิตในความดูแล และผู้สอนดูรายวิชาที่รับผิดชอบ หัวหน้าภาควิชาและผู้ดูแลระบบดูภาพรวมในขอบเขตภาควิชา ส่วนรองคณบดีดูภาพรวมระดับคณะ ผู้ที่ได้รับสิทธิ์อ่านรายงานอย่างเดียวไม่จำเป็นต้องมีสิทธิ์พิจารณาใบลา
รายงานกรองตามปีการศึกษา ภาคการศึกษา ประเภทลา และมิติที่เกี่ยวข้อง เพื่อช่วยติดตามแนวโน้มและเปรียบเทียบข้อมูลภายใต้ขอบเขตเดียวกัน ระบบแสดงเวลาปรับปรุงรายงานล่าสุดให้ผู้ใช้ประเมินความสดของข้อมูลก่อนนำไปประกอบการตัดสินใจ
\subsection{โครงสร้างคลังข้อมูลและนิยามตัวชี้วัด}
คลังข้อมูลแยกข้อมูลวิเคราะห์ออกจากฐานที่ใช้ยื่นและพิจารณาใบลา โดยจัดเก็บสามระดับ ได้แก่ คำขอ รายวิชาที่ได้รับผลกระทบ และขั้นพิจารณา พร้อมข้อมูลนิสิต บุคลากร ภาคการศึกษา และหน่วยงานสำหรับจำแนกผลรายงาน \cite{warehouse}
จำนวนใบลาและวันลารวมคำนวณที่ระดับคำขอ เช่น ใบลาหนึ่งใบจำนวน 3 วันกระทบ 2 วิชา ยังคงนับเป็น 1 ใบและ 3 วัน ส่วนสถิติรายวิชาใช้วันหรือคาบเรียนที่ได้รับผลกระทบ การแยกระดับข้อมูลช่วยป้องกันยอดรวมเพิ่มซ้ำเมื่อใบลาเกี่ยวข้องกับหลายวิชาหรือมีผู้สอนร่วม
ค่าเฉลี่ยข้ามภาคการศึกษารวมเทอมที่ไม่มีใบลาเป็นศูนย์ด้วย การจัดอันดับช่วยระบุวิชาที่มียอดการลาสูงภายใต้ช่วงที่เลือก ส่วนแผนภาพกล่องของผู้สอนแสดงการกระจายยอดการลาของแต่ละวิชาข้ามเทอม ช่วยเปรียบเทียบค่ากลางและความผันผวน โดยยังไม่อธิบายสาเหตุหรือผลต่อคะแนนเรียน
\subsection{กระบวนการ ETL และความสดของรายงาน}
กระบวนการจัดเตรียมรายงานเริ่มจากรวบรวมคำขอ ประวัติพิจารณา และข้อมูลอ้างอิง จากนั้นตรวจความครบถ้วน เชื่อมนิสิต วิชา ผู้สอน และภาคการศึกษา ก่อนนำเข้าสู่คลังข้อมูล การโหลดซ้ำต้องปรับปรุงรายการเดิมโดยไม่เพิ่มจำนวนใบลา และเมื่อพบรายการที่หายจากต้นทางต้องปรับสถานะไม่ให้นำมารวมยอดต่อ
ระบบตรวจคุณภาพข้อมูลก่อนเผยแพร่ผล หากการโหลดล้มเหลวจะยกเลิกการเปลี่ยนแปลงในรอบนั้นและคงรายงานจากรอบที่สำเร็จล่าสุดไว้ การมีประวัติรอบโหลดช่วยแยกปัญหาข้อมูลยังไม่ปรับปรุงออกจากปัญหาการคำนวณรายงาน
ก่อนส่งผลรายงาน ระบบตรวจตัวตน บทบาท และความสัมพันธ์ปัจจุบันกับนิสิตหรือรายวิชาอีกครั้ง ข้อมูลที่เก็บเพื่อวิเคราะห์ย้อนหลังจึงไม่ใช้แทนสิทธิ์เข้าถึงปัจจุบัน แนวทางนี้ช่วยให้รายงานรักษาบริบทประวัติพร้อมจำกัดข้อมูลตามหน้าที่ของผู้ใช้งาน
